\section*{Chapitre \ref{ch:b2:cell-differentiation} --- La différenciation cellulaire~: la cellule musculaire squelettique}

\begin{solution}{exo:b2:cell-differentiation:1}
\omterm{def:b2:cell-differentiation:differentiation}{Différenciation}~: l'expression stable des gènes et des fonctions d'un
type cellulaire. \omterm{def:b2:cell-differentiation:differentiation}{Détermination}~: l'engagement vers un destin avant qu'il
ne s'exprime. \omterm{def:b2:cell-differentiation:differentiation}{Potentialité}~: l'éventail des destins encore ouverts.
Zygote~: totipotent~; cellule de la \omterm{def:b2:mammal-reproduction:implantation}{masse cellulaire interne}~:
pluripotente~; cellule satellite~: unipotente (muscle seulement) --- une
\omterm{def:b2:cell-differentiation:differentiation}{cellule souche} néanmoins~; \omterm{def:b2:cell-differentiation:myogenesis}{fibre musculaire}~: différenciée,
post-mitotique, sans \omterm{def:b2:cell-differentiation:differentiation}{potentialité}.
\end{solution}

\begin{solution}{exo:b2:cell-differentiation:2}
Cellule du dermomyotome (Pax3) $\to$ \omterm{def:b2:cell-differentiation:myogenesis}{myoblaste} induit (\omterm{def:b2:cell-differentiation:myogenesis}{MyoD}, Myf5), qui
se divise tant que dure le FGF $\to$ sortie du cycle, myogénine $\to$
alignement et fusion en un \omterm{def:b2:cell-differentiation:myogenesis}{myotube} $\to$ activation des gènes
musculaires, assemblage des \omterm{def:b2:cell-differentiation:fibre}{myofibrilles} (MRF4) $\to$ fibre mûre à noyaux
périphériques, accompagnée de cellules satellites quiescentes (Pax7).
\end{solution}

\begin{solution}{exo:b2:cell-differentiation:3}
Que le noyau d'une cellule différenciée conserve tous les gènes
nécessaires pour construire un animal entier, et que le cytoplasme de
l'œuf peut réinitialiser son programme~: la \omterm{def:b2:cell-differentiation:differentiation}{différenciation} est
réversible et n'est pas une perte d'ADN. L'expérience n'a pas montré que
tout noyau différencié peut être réinitialisé (la plupart des transferts
ont échoué), ni comment se fait la réinitialisation, ni que la cellule
différenciée elle-même pourrait changer de destin sur place.
\end{solution}

\begin{solution}{exo:b2:cell-differentiation:4}
Les stries Z qui délimitent le \omterm{def:b2:cell-differentiation:fibre}{sarcomère}~; les filaments fins d'actine
qui y sont ancrés~; les filaments épais de myosine centrés sur la ligne
M~; la bande A (filaments épais) et la bande I (filaments fins seuls).
Lors du raccourcissement, les filaments fins glissent vers la ligne M et
les stries Z se rapprochent~: la bande I et la zone H rétrécissent, la
bande A --- la longueur des filaments épais --- ne change pas, et aucun
filament ne raccourcit.
\end{solution}

\begin{solution}{exo:b2:cell-differentiation:5}
$0.2/2.5\times 10^{-6} = \num{80000}$ \omterm{def:b2:cell-differentiation:fibre}{sarcomères}~; $0.2/25\times 10^{-6}
= \num{8000}$ noyaux~; un \omterm{def:b2:cell-differentiation:myogenesis}{myoblaste} par noyau, soit environ \num{8000}
\omterm{def:b2:cell-differentiation:myogenesis}{myoblastes} fusionnés (les cellules satellites en ajoutent d'autres par
la suite).
\end{solution}

\begin{solution}{exo:b2:cell-differentiation:6}
Le FGF induit Id, qui se fixe sur la protéine partenaire dont \omterm{def:b2:cell-differentiation:myogenesis}{MyoD} a
besoin pour former un dimère capable de se lier à l'ADN~; \omterm{def:b2:cell-differentiation:myogenesis}{MyoD} est
présent mais inactif, et les cellules continuent de se diviser. Lorsque
le FGF est retiré, Id se dégrade, \omterm{def:b2:cell-differentiation:myogenesis}{MyoD} se dimérise, la myogénine est
activée, et les cellules sortent du cycle et fusionnent. Des \omterm{def:b2:cell-differentiation:myogenesis}{myoblastes}
produisant Id de façon constitutive ne fusionnent jamais~: ils se
divisent indéfiniment et ne forment pas de muscle.
\end{solution}

\begin{solution}{exo:b2:cell-differentiation:7}
Condition~: $0.4m = m^{2}/(1 + m^{2}) + 0.02$. Pour $m = 2.1$~: membre
de gauche $0.84$, membre de droite $4.41/5.41 + 0.02 = 0.835$~: c'est un
état stationnaire. État bas~: pour $m$ petit, $0.4m \approx m^{2} + 0.02$,
d'où $m \approx
0.055$ (gauche $0.022$, droite $0.023$). Portée à 1, valeur supérieure
au seuil voisin de $0.41$, la cellule monte jusqu'à $2.1$ et y reste~:
elle est engagée.
\end{solution}

\begin{solution}{exo:b2:cell-differentiation:8}
Avec $k = 0.6$~: pour $m = 1$, la dégradation $0.6$ dépasse la synthèse
$0.52$~; pour $m = 0.5$, $0.30 > 0.22$~; pour $m = 2$, $1.2 > 0.82$~: la
droite est partout au-dessus de la courbe au-delà de l'état bas, qui est
le seul état stationnaire. Une cellule dont \omterm{def:b2:cell-differentiation:myogenesis}{MyoD} est dégradé trop vite ne
peut pas maintenir l'état allumé~: elle retombe dans l'état
indifférencié quelle que soit l'\omterm{def:b2:vertebrate-development:induction}{induction} --- pas de muscle.
\end{solution}

\begin{solution}{exo:b2:cell-differentiation:9}
(a) Ni \omterm{def:b2:cell-differentiation:myogenesis}{myoblastes} ni muscle squelettique~: les deux facteurs sont
redondants pour la \omterm{def:b2:cell-differentiation:differentiation}{détermination}, et il faut donc supprimer les deux.
(b) Des \omterm{def:b2:cell-differentiation:myogenesis}{myoblastes} se forment et s'alignent mais ne fusionnent jamais ni
ne fabriquent de protéines musculaires~: la myogénine est nécessaire à la
\omterm{def:b2:cell-differentiation:differentiation}{différenciation}. (c) Presque normale, car Myf5 remplace \omterm{def:b2:cell-differentiation:myogenesis}{MyoD} pour la
\omterm{def:b2:cell-differentiation:differentiation}{détermination}.
\end{solution}

\begin{solution}{exo:b2:cell-differentiation:10}
Dans un fibroblaste, les gènes musculaires se trouvent dans une
chromatine que \omterm{def:b2:cell-differentiation:myogenesis}{MyoD} peut ouvrir~; dans une cellule hépatique, ils sont
peut-être compactés dans de l'hétérochromatine, ou réprimés par les
\omterm{prop:b2:cell-differentiation:transcriptional}{régulateurs maîtres} propres au foie~: la cellule n'est pas compétente.
Test~: traiter les cellules hépatiques par un médicament qui relâche la
chromatine (un inhibiteur des histone désacétylases) avant d'ajouter
\omterm{def:b2:cell-differentiation:myogenesis}{MyoD}, ou ajouter \omterm{def:b2:cell-differentiation:myogenesis}{MyoD} en même temps qu'un facteur qui supprime le
programme hépatique, et mesurer l'expression de la myosine.
\end{solution}

\begin{solution}{exo:b2:cell-differentiation:11}
Une décharge soutenue à basse fréquence élève longtemps le calcium
intracellulaire~; des voies activées par le calcium allument les gènes
de la myosine lente, de la biogenèse mitochondriale, de la myoglobine et
de la croissance des capillaires, et éteignent les isoformes rapides. Le
programme principal de la fibre (état de \omterm{def:b2:cell-differentiation:myogenesis}{MyoD}, architecture des
\omterm{def:b2:cell-differentiation:fibre}{sarcomères}) n'est pas touché~: l'entraînement change les gènes
musculaires exprimés, non l'identité de la cellule, et aucun signal
nerveux ne peut rouvrir le programme du cartilage, fermé lors de la
\omterm{def:b2:cell-differentiation:differentiation}{détermination}.
\end{solution}

\begin{solution}{exo:b2:cell-differentiation:12}
L'état allumé d'un facteur auto-activateur est un point stable d'une
boucle de rétroaction, maintenu par une synthèse continue, et non par
une modification des gènes~; les marques de la chromatine ajoutent une
seconde couche qui garde fermés les gènes éteints. La réinitialisation
exige donc de ramener le circuit sous son seuil et de rouvrir la
chromatine en même temps, ce que le cytoplasme de l'œuf ou quatre
facteurs ne réussissent que dans une faible fraction des cellules~:
c'est possible, parce que rien n'est perdu, mais peu efficace, parce
que deux verrous indépendants doivent être crochetés ensemble.
\end{solution}

\begin{solution}{pb:b2:cell-differentiation:1}
\textbf{1.} $0.2/2.5\times 10^{-6} = \num{80000}$ \omterm{def:b2:cell-differentiation:fibre}{sarcomères}~;
$0.2/25\times 10^{-6} = \num{8000}$ noyaux.
\textbf{2.} Environ \num{8000} par fibre~; $4\times 10^{9}$ pour le
muscle.
\textbf{3.} $\pi(30\times 10^{-6})^{2}\times 0.2 = \qty{5.7e-10}{m^3}$
par fibre~; $2.8\times 10^{-4}$ \unit{m^3}, soit \qty{0.28}{L}, pour le
muscle.
\textbf{4.} $5.7\times 10^{-10}/8000 = \qty{7e-14}{m^3} =
\qty{70000}{\micro m^3}$~: l'équivalent de 35 cellules ordinaires par
noyau.
\textbf{5.} L'intérieur est rempli de \omterm{def:b2:cell-differentiation:fibre}{myofibrilles} qui doivent courir
sans interruption d'un bout à l'autre~; des noyaux situés en bordure
laissent le réseau continu et se trouvent près de la membrane et de ses
signaux.
\textbf{6.} $4\times 10^{9}/2000 = 2\times 10^{6}$~: 21 doublements,
10.5 jours.
\textbf{7.} $0.4m = m^{2}/(1 + m^{2}) + 0.02$. Pour $m = 0.055$~:
gauche $0.022$, droite $0.003 + 0.02 = 0.023$.
\textbf{8.} Pour $m = 2.1$~: gauche $0.84$, droite $0.815 + 0.02 = 0.835$.
\textbf{9.} Pour $m = 0.41$~: gauche $0.164$, droite $0.144 + 0.02 = 0.164$.
Juste en dessous, la synthèse est inférieure à la dégradation, et $m$
descend vers $0.055$~; juste au-dessus, la synthèse dépasse la
dégradation et $m$ monte vers $2.1$~: tout écart s'amplifie --- l'état
est instable.
\textbf{10.} $0.6 > 0.41$~: la première cellule passe à l'état allumé et
se différencie~; $0.3 < 0.41$~: la seconde revient à $0.055$ et ne se
différencie pas.
\textbf{11.} Avec $b = 0.4$, la courbe de synthèse part de $0.4$ et
reste au-dessus de la droite $0.4m$ jusqu'à $m \approx 3.3$~: le
croisement bas a disparu. La cellule se différencie spontanément, sans
\omterm{def:b2:vertebrate-development:induction}{induction}.
\textbf{12.} La division divise $m$ par deux~; les cellules filles
héritent de $2.1/2 =
1.05$, au-dessus du seuil $0.41$, et remontent à $2.1$~: l'état est
rétabli dans chaque cellule fille. L'engagement persiste tant que le
niveau réduit de moitié reste au-dessus du seuil.
\textbf{13.} $4$ par mm $\times$ \qty{200}{mm} $= 800$ cellules
satellites par fibre~; \qty{10}{\%} des \num{8000} noyaux, soit
\num{800}, doivent être remplacés.
\textbf{14.} Le segment lésé contient $80$ cellules satellites~;
$800/80 =
10$, donc 4 doublements ($\times 16$, soit \num{1280}~: 800 fusionnent,
480 restent pour reconstituer le stock) --- \qty{72}{h}, trois jours.
\textbf{15.} Trois jours de divisions contre deux à trois semaines
observées~: le reste correspond à l'élimination des débris par les
macrophages, au délai d'activation, à la fusion, à l'assemblage des
\omterm{def:b2:cell-differentiation:fibre}{myofibrilles}, à la réinnervation et à la maturation du nouveau segment.
\textbf{16.} \omterm{def:b2:cell-differentiation:myogenesis}{MyoD} (et la myogénine) doivent réactiver les gènes
musculaires dans les cellules activées~; division et \omterm{def:b2:cell-differentiation:differentiation}{différenciation}
s'excluent --- Id doit diminuer et le cycle s'arrêter avant que la
myogénine puisse agir et que la fusion ait lieu.
\textbf{17.} $0.95^{n} = 0.5$~: $n = 13.5$, soit environ quatorze
lésions.
\textbf{18.} $800\times 0.95^{20} = 800\times 0.36 \approx 290$.
\textbf{19.} Chaque contraction déchire des fibres, chaque déchirure
sollicite le stock, et le stock diminue de quelques pour cent à chaque
fois~; au bout d'années, il est épuisé, les fibres déchirées ne sont
plus reconstruites, et les fibroblastes comblent les vides avec du
collagène.
\textbf{20.} $(80/60)^{2} = 1.78$~; noyaux $8000\times 1.78 = \num{14200}$~:
environ \num{6200} supplémentaires fournis par les cellules satellites.
\textbf{21.} $0.28\times 1.78 = \qty{0.5}{L}$.
\textbf{22.} Ce qui a changé~: les gènes des isoformes de la myosine, la
densité des mitochondries et des capillaires, la myoglobine. Ce qui n'a
pas changé~: l'identité musculaire (état de \omterm{def:b2:cell-differentiation:myogenesis}{MyoD}), l'architecture des
\omterm{def:b2:cell-differentiation:fibre}{sarcomères}, les noyaux et leur génome.
\textbf{23.} Les fibres sont post-mitotiques et ne peuvent pas se
multiplier~; les seules voies sont des fibres plus épaisses et davantage
de noyaux par fibre, fournis par les cellules satellites.
\textbf{24.} Entraînement~: un épaississement limité, puisque la fibre
ne peut pas ajouter de noyaux~; lésion~: pas de réparation --- les
segments détruits sont remplacés par du tissu fibreux, comme dans la
myopathie.
\textbf{25.} \num{8000} noyaux par fibre~; seuil $m \approx 0.41$~;
environ trois jours de divisions des cellules satellites pour une
réparation.
\end{solution}
